%% ============================================================
%%  AI, Big Data, and Sierra Leone's Energy Transition:
%%  A Data-to-Decision Framework for Utility Reform
%%  and Distributed Electrification
%%
%%  Thoronka & Alpha (2025)  —  Cities: Smarter Futures
%%  Compile:  pdflatex → bibtex → pdflatex → pdflatex
%%  All .png figures must be in the same folder as this .tex
%% ============================================================

\documentclass[preprint,12pt]{elsarticle}

%% ── Encoding & typography ────────────────────────────────────
\usepackage[T1]{fontenc}
\usepackage[utf8]{inputenc}

%% ── Mathematics ──────────────────────────────────────────────
\usepackage{amsmath,amssymb,mathtools}

%% ── Layout ───────────────────────────────────────────────────
\usepackage{geometry}
\geometry{a4paper, top=2.8cm, bottom=2.8cm, left=2.6cm, right=2.6cm}
\usepackage{float, caption, subcaption}
\usepackage{setspace, parskip}

%% ── Tables ───────────────────────────────────────────────────
\usepackage{booktabs, tabularx, array, multirow, colortbl, makecell}

%% ── Graphics ─────────────────────────────────────────────────
\usepackage{graphicx}

%% ── Colour system ────────────────────────────────────────────
\usepackage{xcolor}
\definecolor{hdrblue}{HTML}{1E3A5F}
\definecolor{hdrgreen}{HTML}{064E3B}
\definecolor{hdramber}{HTML}{451A03}
\definecolor{hdrpurple}{HTML}{2E1065}
\definecolor{hdrteal}{HTML}{134E4A}
\definecolor{rowA}{HTML}{EFF6FF}
\definecolor{rowB}{HTML}{ECFDF5}
\definecolor{rowC}{HTML}{FFFBEB}
\definecolor{rowD}{HTML}{F5F3FF}
\definecolor{midGray}{HTML}{E5E7EB}
\definecolor{lightGray}{HTML}{F9FAFB}
\definecolor{critRed}{HTML}{FEE2E2}
\definecolor{critBdr}{HTML}{CB181D}
\definecolor{passGrn}{HTML}{DCFCE7}
\definecolor{passBdr}{HTML}{15803D}
\definecolor{loopOrange}{HTML}{EA580C}
\definecolor{loopRed}{HTML}{991B1B}
\definecolor{loopBg}{HTML}{FFF7ED}
\definecolor{boxDRI}{HTML}{EFF6FF}
\definecolor{boxDBI}{HTML}{FFF7ED}
\definecolor{boxDP}{HTML}{F5F3FF}
\definecolor{boxMC}{HTML}{ECFDF5}

%% ── TColorBox ────────────────────────────────────────────────
\usepackage{tcolorbox}
\tcbuselibrary{skins, breakable, theorems}

%% Feedback loop box (the centrepiece visual)
\tcbset{
  loopbox/.style={
    enhanced, breakable,
    colback=loopBg,
    colframe=loopRed,
    boxrule=1.8pt,
    arc=6pt,
    left=12pt, right=12pt, top=10pt, bottom=10pt,
    drop shadow={opacity=0.18},
  }
}

%% Metric box (numbered equations with context)
%% Named metric box styles (avoid ## parameter issues)
\tcbset{
  metDRI/.style={enhanced, breakable,
    colback=boxDRI, colframe=hdrblue, boxrule=1.2pt, arc=4pt,
    fonttitle=\bfseries\small\color{white},
    attach boxed title to top left={yshift=-2.5mm, xshift=6mm},
    boxed title style={colback=hdrblue, rounded corners, boxrule=0pt},
    left=8pt, right=8pt, top=10pt, bottom=8pt,
    drop shadow={opacity=0.10}},
  metDBI/.style={enhanced, breakable,
    colback=boxDBI, colframe=loopOrange, boxrule=1.2pt, arc=4pt,
    fonttitle=\bfseries\small\color{white},
    attach boxed title to top left={yshift=-2.5mm, xshift=6mm},
    boxed title style={colback=loopOrange, rounded corners, boxrule=0pt},
    left=8pt, right=8pt, top=10pt, bottom=8pt,
    drop shadow={opacity=0.10}},
  metDP/.style={enhanced, breakable,
    colback=boxDP, colframe=hdrpurple, boxrule=1.2pt, arc=4pt,
    fonttitle=\bfseries\small\color{white},
    attach boxed title to top left={yshift=-2.5mm, xshift=6mm},
    boxed title style={colback=hdrpurple, rounded corners, boxrule=0pt},
    left=8pt, right=8pt, top=10pt, bottom=8pt,
    drop shadow={opacity=0.10}},
  metMC/.style={enhanced, breakable,
    colback=boxMC, colframe=hdrgreen, boxrule=1.2pt, arc=4pt,
    fonttitle=\bfseries\small\color{white},
    attach boxed title to top left={yshift=-2.5mm, xshift=6mm},
    boxed title style={colback=hdrgreen, rounded corners, boxrule=0pt},
    left=8pt, right=8pt, top=10pt, bottom=8pt,
    drop shadow={opacity=0.10}},
}

%% Key finding box
\tcbset{
  findingbox/.style={
    enhanced,
    colback=rowD,
    colframe=hdrpurple,
    fonttitle=\bfseries\small\color{white},
    attach boxed title to top center={yshift=-2.5mm},
    boxed title style={colback=hdrpurple, rounded corners},
    arc=4pt, left=8pt, right=8pt, top=10pt, bottom=8pt,
    boxrule=1.2pt,
  }
}

%% ── Links ────────────────────────────────────────────────────
\usepackage[unicode=true, hidelinks]{hyperref}

%% ── Lists ────────────────────────────────────────────────────
\usepackage{enumitem}

%% ── Captions ─────────────────────────────────────────────────
\captionsetup{font=small, labelfont=bf, justification=justified,
              singlelinecheck=false, skip=5pt}

%% ── Table helpers ────────────────────────────────────────────
\newcommand{\Th}[1]{\cellcolor{hdrblue}\color{white}\bfseries\small #1}
\newcommand{\ThG}[1]{\cellcolor{hdrgreen}\color{white}\bfseries\small #1}
\newcommand{\ThA}[1]{\cellcolor{hdramber}\color{white}\bfseries\small #1}
\newcommand{\ThP}[1]{\cellcolor{hdrpurple}\color{white}\bfseries\small #1}
\newcommand{\ThT}[1]{\cellcolor{hdrteal}\color{white}\bfseries\small #1}
\newcommand{\Fail}[1]{\cellcolor{critRed}\textcolor{critBdr}{\bfseries #1}}
\newcommand{\Pass}[1]{\cellcolor{passGrn}\textcolor{passBdr}{\bfseries #1}}

\journal{Cities: Smarter Futures}

%% ─────────────────────────────────────────────────────────────
\begin{document}

\begin{frontmatter}

\title{AI, Big Data, and Sierra Leone's Energy Transition:
       A Data-to-Decision Framework for Utility Reform
       and Distributed Electrification}

\author[tu]{Jeremiah Thoronka\corref{cor1}}
\ead{jeremiahthoronka1@gmail.com}
\cortext[cor1]{Corresponding author}

\author[cmu]{Mohamed Alpha}
\ead{malpha@andrew.cmu.edu}

\address[tu]{School of Global Studies, Thammasat University,
             Bangkok, Thailand}
\address[cmu]{Carnegie Mellon University Africa, Kigali, Rwanda}

\begin{abstract}
Sierra Leone's energy transition is conventionally narrated through
the twin deficits of inadequate generation capacity and incomplete
network extension. This article contends that a third, deeper deficit
has been systematically overlooked: a data-to-decision gap that
renders the sector incapable of observing, interpreting, and acting
upon its own operational conditions---even as the electricity
landscape becomes more distributed, heterogeneous, and data-intensive.
Drawing on sociotechnical transition theory, energy justice
scholarship, and algorithmic governance literature, we develop the
\textbf{AI--Big Data Catalyst Model for Sierra Leone (AIBD-Cat-SL)}:
a four-layer architecture in which big data produces observability
and AI produces actionability, but where governance---not
algorithmic sophistication---determines whether digitalisation serves
all communities or only those already data-visible.

Two original quantitative contributions operationalise the framework
on real published data. The \textbf{Double Blindspot Index (DBI)},
a weighted composite of access deprivation, poverty burden, and data
opacity, reveals that electricity deprivation and data invisibility
co-locate across Sierra Leone's 14 administrative districts with
$r = 0.972$ (DBI range: 15.8 in Freetown to 88.6 in Falaba). The
\textbf{first demographic parity audit of AI-driven electrification
planning in any low- and middle-income country (LMIC)} demonstrates
that a profit-optimised allocation AI produces a parity gap of
$\Delta_{\mathrm{DP}} = 0.406$, assigning zero investment probability
to every high-poverty district. A single Layer~4 governance rule
reduces this gap by \textbf{92\%} to $\Delta_{\mathrm{DP}} = 0.034$
without reducing efficiency. Monte Carlo validation across 15,000
simulations confirms both findings are structural properties of
Sierra Leone's data environment, not artefacts of specific
parameter choices. Eight testable propositions, a phased
implementation roadmap, and a concrete procurement precedent
translate the framework into actionable guidance.
\end{abstract}

\begin{keyword}
artificial intelligence \sep big data \sep energy transition \sep
Sierra Leone \sep electrification \sep utility reform \sep
energy justice \sep predictive maintenance \sep non-technical
losses \sep digital governance \sep demographic parity \sep
data readiness
\end{keyword}

\end{frontmatter}

\onehalfspacing

%% =============================================================
\section{Introduction}
\label{sec:intro}
%% =============================================================

Ghana's Electricity Company (ECG) reduced technical and commercial
losses from nearly 27\% to 18\% between 2012 and 2019. It did not
build new power plants to do it. It installed prepaid meters,
deployed feeder-level monitoring, and used the resulting data to
identify exactly where electricity was disappearing. The physical
infrastructure barely changed. The information environment changed
completely. A decade later, Sierra Leone's EDSA still reports
T\&D losses of approximately 30\%---ten percentage points above
the global benchmark and twice the level of a well-governed
African utility. The gap between Ghana and Sierra Leone in loss
performance is not primarily an engineering gap. It is an
information governance gap.

Sierra Leone's energy transition is entering a moment of unusual
strategic significance. In recent years, the country has witnessed
an intensification of policy ambition around electrification,
renewable energy deployment, utility reform, and institutional
strengthening. Yet the structural conditions of the sector remain
deeply challenging. Persistent energy poverty, uneven service
reliability, high technical and commercial losses, and fragile
utility performance continue to undermine the developmental promise
of electrification. These constraints are typically interpreted
through inadequate generation capacity, weak transmission and
distribution infrastructure, and limited network reach. While
such factors remain critically important, they do not fully capture
the deeper operational and institutional dynamics at play.

This paper argues that Sierra Leone's transition must be
understood not only as a problem of infrastructural insufficiency,
but also as a problem of \textit{informational weakness}: a
system-level inability to observe, interpret, and act upon
rapidly changing technical, spatial, and socio-economic conditions
across a diversifying electricity landscape. This argument is
increasingly salient because the underlying architecture of
electrification in Sierra Leone is no longer singular, linear,
or purely grid-centric. The country's transition is progressively
assuming the characteristics of a \textit{hybrid energy assemblage}
in which central grid expansion coexists with mini-grids, solar
home systems, hybrid power configurations, imported electricity,
and emerging renewable resources. Such pluralisation fundamentally
transforms the epistemic and operational demands placed upon the
energy system.

At the heart of this paper is an uncomfortable discovery.
When we ran a cost-optimised AI allocation model on Sierra Leone's
actual district data, the result was clear and immediate: every
single high-poverty district received an investment probability
of exactly zero. Not 5\%. Not 2\%. Zero. The AI was not making
a mistake. It was doing what cost-optimisation algorithms do:
favouring geography that is already close to the existing grid.
But that geography is the same geography that received investment
over the past fifty years of unequal infrastructure development.
The AI inherited the past and projected it as the future.

This is the problem the paper addresses---and solves. The solution
is not a better algorithm. It is a governance rule. One rule.

A growing body of scholarship confirms that AI and big data
analytics can enhance load forecasting, predictive maintenance,
fault detection, revenue protection, and infrastructure planning
\cite{Ahmad2021, Arevalo2024, Chatterjee2022, Yousef2023}. Yet
much of this literature either speaks in universal terms or
assumes higher-capacity institutional environments where data
infrastructures are more mature and governance institutions are
better equipped. Far less attention has been devoted to how
digital intelligence should be sequenced and governed in
low-income settings where the central challenge is not optimisation
within already legible systems, but the very \textit{production
of legibility itself}.

Against this background, this paper develops the
\textbf{AI--Big Data Catalyst Model for Sierra Leone (AIBD-Cat-SL)},
a four-layer conceptual architecture that links data infrastructure,
analytical intelligence, decision routines, and the enabling
environment of governance and public accountability. The article
advances four interrelated claims: (1)~the primary value of AI
in capacity-constrained settings is augmenting human
decision-making, not automating it away; (2)~the benefits of
AI are conditional on data quality, interoperability, and
governance; (3)~digitalisation strengthens utility viability
and electrification strategy only when designed around public
value, not just cost efficiency; and (4)~Sierra Leone is an
important theoretical site for studying digital energy transitions
precisely because it reveals the full tension between digital
ambition and institutional fragility.

The paper makes four principal contributions. First, it
reconceptualises Sierra Leone's energy transition as a
\textit{three-deficit} challenge---adding a data-to-decision
deficit to the familiar infrastructural and financial ones.
Second, it introduces the AIBD-Cat-SL framework as a structured
architecture for how observability, analytics, institutional
action, and governance must align. Third, it delivers the
Double Blindspot Index and the first LMIC algorithmic parity
audit as original, reproducible, open-data contributions.
Fourth, it validates these findings through 15,000 Monte Carlo
simulations and translates them into a phased implementation
roadmap with specific procurement requirements.

%% =============================================================
\section{Methodology}
\label{sec:method}
%% =============================================================

This study adopts a \textbf{theory-building with analytical
demonstration} design. The primary contribution is conceptual---
the AIBD-Cat-SL framework and eight testable propositions---
situated within an interpretive-analytical epistemology that
treats energy transitions as sociotechnical transformations rather
than purely technical events. Three analytical demonstrations
then operationalise the framework's claims using exclusively real,
open-access data. Monte Carlo simulation validates the structural
robustness of the findings. Table~\ref{tab:data} lists all
data sources.

\begin{table}[H]
\centering
\caption{\textbf{Data sources, provenance, and figure linkage.}
All sources are open-access and independently verifiable.}
\label{tab:data}
\setlength\extrarowheight{3pt}
\begin{tabularx}{\textwidth}{|>{\arraybackslash}p{3.4cm}
                              |>{\arraybackslash}X
                              |>{\centering\arraybackslash}p{1.0cm}
                              |>{\centering\arraybackslash}p{1.6cm}|}
\hline
\Th{Source} & \Th{Coverage and purpose} & \Th{Scale} & \Th{Figures}\\
\hline
\rowcolor{rowA}
World Bank WDI \cite{WorldBankWDI} &
  National electrification and T\&D losses 1990--2023 &
  Nat. & 1, 4 \\
\hline
GADM v4.1 \cite{GADM} &
  SL administrative boundaries, 14 districts &
  Dist. & 2, 3 \\
\hline
\rowcolor{rowA}
ESMAP/MTENR 2023 \cite{ESMAP2023} &
  District access rates, NTL and loss diagnostics &
  Dist. & 2, 3, 6, 8 \\
\hline
WB Poverty Assessment SL \cite{WorldBankWDI} &
  District poverty headcount (\$2.15/day PPP) &
  Dist. & 2, 7 \\
\hline
\rowcolor{rowA}
DHS Programme SL 2019 \cite{DHS2019} &
  Household energy access, gender disaggregation &
  HH & 2, 3 \\
\hline
Statistics Sierra Leone \cite{StatsSL2021} &
  District population and settlement distribution &
  Dist. & 2, 7 \\
\hline
\rowcolor{rowA}
UNDP HDR SL 2022 \cite{UNDP2022} &
  District Human Development Index estimates &
  Dist. & 2 \\
\hline
SEforALL GEP \cite{Korkovelos2019} &
  Technology-matching rules for electrification &
  Dist. & 5 \\
\hline
\rowcolor{rowA}
EDSA Annual Report \cite{EDSA2022} &
  Data legibility dimensions for DRI construction &
  Utility & 8 \\
\hline
AfDB 2022 \cite{AfDB2022} &
  NTL rate calibration (15--25\% of total losses) &
  Nat. & 6 \\
\hline
\end{tabularx}
\end{table}

The analytical strategy unfolds through three stages.
\textit{Integrative synthesis}: six bodies of scholarship are
brought into dialogue---sociotechnical transitions, energy justice,
data infrastructure studies, AI applications in energy systems,
electrification planning, and AI governance. \textit{Contextual
embedding}: concepts from those literatures are interpreted through
Sierra Leone's distinctive conditions---persistent utility fragility,
high losses, weak observability, hybrid service modalities, and
constrained institutional capacity. This stage is not merely
illustrative; it is how the paper avoids uncritically transferring
models from higher-capacity environments. \textit{Conceptual model
construction}: patterns identified across the literature are
translated into the AIBD-Cat-SL layered architecture, with
propositions formulated to guide subsequent empirical work.
The Monte Carlo validation in Section~\ref{sec:mc} then tests
whether the framework's empirical claims are structurally robust.

%% =============================================================
\section{The Data-to-Decision Challenge}
\label{sec:challenge}
%% =============================================================

\subsection{The Feedback Loop That Makes Everything Worse}

Before introducing the framework, it helps to name the central
structural problem as precisely as possible---because understanding
why it exists makes the solution obvious. The problem is a
self-reinforcing feedback loop that connects physical poverty,
data poverty, and algorithmic exclusion into a single,
mutually reinforcing system. It works like this:

\begin{tcolorbox}[loopbox]
\centering\large
\textbf{\textcolor{loopRed}{No electricity}}
\hspace{0.5em}
\textcolor{loopOrange}{$\boldsymbol{\longrightarrow}$}
\hspace{0.5em}
\textbf{\textcolor{loopRed}{No data}}
\hspace{0.5em}
\textcolor{loopOrange}{$\boldsymbol{\longrightarrow}$}
\hspace{0.5em}
\textbf{\textcolor{loopRed}{AI ignores you}}
\hspace{0.5em}
\textcolor{loopOrange}{$\boldsymbol{\longrightarrow}$}
\hspace{0.5em}
\textbf{\textcolor{loopRed}{Still no electricity}}

\medskip
{\small\itshape\textcolor{hdramber}{
Under-investment produces energy poverty. Energy poverty
produces data scarcity. Data scarcity produces AI systems that
cannot see deprived communities. AI systems that cannot see
deprived communities recommend investment elsewhere.
Elsewhere gets more electricity, more data, and better AI.
The cycle deepens with every round.}}
\end{tcolorbox}

\medskip
This is not a theoretical possibility. It is what the data
shows. Our Double Blindspot Index, introduced formally in
Section~\ref{sec:math}, measures the co-location of energy
deprivation and data invisibility across Sierra Leone's 14
administrative districts. The correlation is $r = 0.972$---
near-perfect, in the direction you would expect from the
feedback loop described above. Breaking this cycle requires
simultaneously addressing physical infrastructure and digital
infrastructure. Doing one without the other simply means the
AI-enabled investments flow to the already-connected, while
the unconnected fall further behind.

\subsection{Electrification Heterogeneity and System Complexity}

The contemporary electrification challenge in Sierra Leone
extends well beyond the numerical expansion of connections.
It increasingly concerns the governance of a rapidly diversifying
energy landscape in which access is mediated through multiple
technological and service modalities. National planning frameworks
increasingly envisage a mixed electrification pathway in which
future access gains will be delivered through continued main-grid
expansion alongside mini-grids and stand-alone solar systems.
This orientation is rational from the standpoint of geography,
dispersed settlement patterns, and least-cost planning logic.
Yet it introduces a fundamental systemic tension: as the physical
architecture of electrification becomes more fragmented and
heterogeneous, the informational and coordinative architecture
of the sector must become correspondingly more integrated.

A hybrid electricity landscape requires coordination across
infrastructures that differ in ownership structure, technical
design, metering arrangements, telemetry capacity, maintenance
cycles, tariff regimes, and consumer interfaces. The central
issue is \textit{observability}: whether planners, utilities,
regulators, and service providers possess the informational depth
necessary to render a dispersed and differentiated system
governable in practice. Figure~\ref{fig:macro} documents the
consequences of weak observability: access has near-quadrupled
since 2000 while T\&D losses have remained stuck near 30\%.

\begin{figure}[H]
\centering
\includegraphics[width=\linewidth]{Fig1_Expansion_Efficiency.png}
\caption{\textbf{Sierra Leone energy system: the
expansion--efficiency paradox, 1990--2023.}
National electricity access grew from 7.7\% to 35.5\% since
2000 while T\&D losses declined only modestly from
$\approx$44\% to 30\%, remaining far above the global
benchmark of $\approx$20\% and the high-income country
norm of $\approx$8\%. Panel~(b) shows the West Africa
peer comparison: Sierra Leone consistently sits at the
bottom of its regional peer group. Ghana at 87\%,
Senegal at 75\%, and even Liberia at 31\% have all
pulled ahead. This is not an infrastructure story alone.
\textit{Source:} World Bank WDI \cite{WorldBankWDI}.}
\label{fig:macro}
\end{figure}

\subsection{Utility Fragility, Losses, and Service Continuity}

High T\&D losses are not simply a technical inefficiency.
They are the symptom of a utility that cannot see itself
clearly enough to act. EDSA lacks reliable feeder-level
visibility, an up-to-date customer database, tamper detection
capacity, synchronised outage records, and integrated billing
intelligence. As a result, it cannot identify where losses
originate, which assets are failure-prone, or which
interventions would yield the greatest financial return.
The value of AI in this context is not primarily automation.
It is the ability to see what was previously invisible---to
make the invisible legible and the legible actionable.

\subsection{Institutional Reform and the Digital Capability Gap}

Sierra Leone's energy agenda is, in principle, increasingly
reform-oriented. Policy documents call for lower losses,
improved reliability, stronger mini-grid markets, and more
rational investment planning. Yet implementation failure
frequently arises not from the absence of policy ambition,
but from the inability to translate reform goals into everyday
operational routines. Data systems are still treated as
ancillary administrative tools rather than constitutive
infrastructure for sector governance. This must change.
In a distributed, data-intensive electricity transition,
digital capability is not an optional add-on. It is the
institutional infrastructure through which reliability,
coordination, and system intelligence become possible.

%% =============================================================
\section{Theoretical Foundations}
\label{sec:theory}
%% =============================================================

Three interconnected theoretical pillars anchor this paper.

\subsection{Sociotechnical Transitions and Epistemic
Reconfiguration}

The multi-level perspective in sociotechnical transition theory
\cite{Geels2002} insists that energy transitions are not just
technology swaps---they are system-wide reconfigurations of
infrastructures, institutions, regulatory arrangements, user
practices, and societal expectations. AI and big data are
not neutral instruments introduced into a stable system. They
are components of a broader \textit{regime reconfiguration}
in which the informational basis of coordination, control, and
adaptation is itself being transformed.

We extend this framework by foregrounding
\textit{epistemic reconfiguration} as a necessary condition
of regime transformation: transitions require not just new
infrastructure but new ways of seeing the system, measuring
its states, and coordinating action across heterogeneous
service modalities. Data standards, telemetry architectures,
and interoperable information flows are not peripheral
complements to transition---they are part of the architecture
through which transition becomes governable at all \cite{Kitchin2014}.

\subsection{Energy Justice and Algorithmic Accountability}

Energy justice scholarship demands that energy systems be
evaluated not only by efficiency metrics but by their
distributive, procedural, and recognitional implications
\cite{Sovacool2015}. Digital energy systems are never socially
neutral. They shape who is seen, whose needs are prioritised,
and whose vulnerabilities are ignored. Decisions generated
or mediated through digital systems---maintenance
prioritisation, disconnection thresholds, subsidy targeting,
infrastructure allocation---carry significant social and
political consequences.

For Sierra Leone, this lens is especially important because
the substantive meaning of energy transition cannot be reduced
to access in the abstract. The question is whether the
transition delivers reliable, affordable, and socially
meaningful services capable of supporting livelihoods,
education, healthcare, and human dignity. Reliability is not
just a technical performance indicator. It is a justice
variable. Unreliable electricity disproportionately harms
those with the fewest coping mechanisms---which in Sierra
Leone means communities in Falaba, Koinadugu, and Pujehun,
not Freetown.

\subsection{Data Infrastructures and Algorithmic Governance}

\citet{Kitchin2014} argues that big data must be understood as
an infrastructural assemblage---constituted through collection,
classification, standardisation, and governance---not merely a
large volume of information. The effectiveness of AI depends
fundamentally on the quality and governance of the data
environment in which it is embedded. In many lower-capacity
settings, the limiting factor is not the absence of algorithms
but the absence of usable data pipelines, interoperable
standards, secure exchange protocols, and the organisational
routines capable of transforming raw information into
operational intelligence.

The literature on algorithmic governance further underscores
that analytics-based decision-making is inseparable from
questions of accountability, transparency, and oversight
\cite{Barocas2019}. These concerns become especially acute
in energy systems where algorithmic outputs influence system
security, household welfare, subsidy allocation, and tariff
enforcement. AI adoption should therefore not be approached
as a software procurement question. It is a process of
institutional embedding in which data, models, human
judgement, organisational workflows, and governance
safeguards must be configured into a coherent actionable whole.

%% =============================================================
\section{Literature Review: Where AI Adds Value,
and Where It Fails}
\label{sec:lit}
%% =============================================================

\subsection{Big Data as the Infrastructure of Observability}

The most consequential contribution of big data to energy
systems lies not in advanced analytics themselves, but in
the \textit{production of enhanced system observability}:
the extent to which system states, asset conditions,
consumption dynamics, fault events, and service disruptions
can be rendered measurable and interpretable across time
and space. In electricity systems, the relevant data
environment encompasses smart meter data, SCADA systems,
feeder and transformer sensors, outage logs, billing records,
mini-grid telemetry, and remote sensing inputs. When
systematically integrated, these streams enable utilities
and planners to move beyond fragmented episodic visibility
toward a coherent understanding of where the system is
underperforming and where intervention would generate the
greatest value.

The electrification planning literature demonstrates that
geospatial tools improve technology-matching decisions, but
their effectiveness depends fundamentally on the granularity,
timeliness, and interoperability of the underlying data
\cite{Korkovelos2019}. In low-access settings, data scarcity
is itself a material planning risk: it misallocates capital,
misjudges lifecycle costs, and reproduces spatial inequities.

\subsection{AI as a Decision-Support Catalyst}

AI in energy systems generates its greatest value in four
areas: \textit{forecasting}, improving projections of demand,
generation, and load variability \cite{Yousef2023};
\textit{predictive maintenance}, identifying vulnerable assets
before breakdown occurs---especially important where maintenance
resources are scarce \cite{Chatterjee2022};
\textit{anomaly detection}, flagging irregular consumption,
theft signatures, and incipient faults that elude manual
analysis \cite{Ahmad2021}; and \textit{optimisation},
coordinating distributed energy resources under conditions
of uncertainty \cite{Arevalo2024}.

These applications generate most value in environments
characterised by high decision complexity, limited response
capacity, and persistent informational asymmetry. This is
precisely the environment Sierra Leone faces. Low-capacity
systems are not unsuitable candidates for AI; they may stand
to benefit most from analytical tools that enhance
prioritisation and the allocation of scarce managerial
attention. Yet this benefit is conditional---and it is that
conditionality that the AIBD-Cat-SL framework is designed to address.

\subsection{Why Adoption Fails}

Many AI initiatives in the energy sector fail to progress
from pilot to sustained practice. Five barriers recur:
data fragmentation that undermines model training; institutional
processes not redesigned to act on analytical insights; low
trust in opaque systems; heightened cybersecurity exposure;
and inadequate workforce readiness \cite{Ahmad2021, Danish2023,
Machlev2022}. These barriers mirror Sierra Leone's structural
conditions almost exactly. The analytical error is to
conceptualise AI adoption as a narrowly technical undertaking.
It is an organisational, institutional, and governance challenge.
Whether AI scales depends less on the algorithm's elegance
than on the sociotechnical arrangements into which it is embedded.

%% =============================================================
\section{The AIBD-Cat-SL Framework}
\label{sec:framework}
%% =============================================================

\subsection{Conceptual Foundation}

The AIBD-Cat-SL rests on a foundational distinction. Big data
constitutes the enabling infrastructure of
\textit{observability}: it creates the ability to see the
system with sufficient precision to govern it. AI constitutes
the catalytic architecture of \textit{actionability}: it
transforms that visibility into prioritised, timely, and
institutionally credible decisions. But neither observability
nor actionability produces equity without explicit governance
design. Figure~\ref{fig:architecture} presents the full
four-layer architecture.

\begin{figure}[H]
\centering
\includegraphics[width=\linewidth]{DiagA_Architecture.png}
\caption{\textbf{AIBD-Cat-SL: four-layer sociotechnical
architecture for Sierra Leone's digital energy transition.}
Layer~1 (blue) produces observability through integrated
data infrastructure. Layer~2 (teal) converts data into
actionable intelligence. Layer~3 (amber) embeds analytics
in organisational routines and consequential decisions.
Layer~4 (purple) is the decisive governance layer---it
determines whether the system serves equity or perpetuates
exclusion. The dashed recursive feedback arc shows how
governance mandates drive data investment, better data
enables more reliable AI, more reliable AI improves
decisions, and improved decisions generate resources for
further governance strengthening.}
\label{fig:architecture}
\end{figure}

\paragraph{Layer~1---Data Foundation.}
The data foundation produces system observability through
integrated metering, feeder monitoring, outage records,
billing databases, mini-grid telemetry, geospatial layers,
and remote sensing. Its function is not to accumulate
information but to establish credible observability through
quality assurance, auditability, and privacy-aware governance.
In Sierra Leone, this layer is the binding constraint:
seven of 14 districts score below 25 on the Data Readiness
Index (Section~\ref{sec:dri}), placing them in a tier
where AI outputs would be unreliable regardless of
algorithmic sophistication.

\paragraph{Layer~2---Intelligence Layer.}
The intelligence layer converts heterogeneous data into
interpretable insights: demand forecasting, predictive
maintenance, anomaly detection, loss analytics, and
equity-constrained planning. AI augments rather than
displaces human judgement. In a resource-constrained utility,
the value of AI is that it helps a limited team focus scarce
attention where it matters most.

\paragraph{Layer~3---Decision and Automation.}
Analytics acquire practical significance only when embedded
in organisational routines: maintenance schedules, inspection
protocols, dispatch decisions, investment planning. The true
threshold of digital transition is crossed not when algorithms
are introduced, but when analytical outputs begin to reshape
organisational behaviour in durable, consistent ways.

\paragraph{Layer~4---Enabling Environment.}
The enabling environment---regulatory support, workforce
capacity, cybersecurity, AI governance, and demographic parity
mandates---determines whether the system produces public value
or perpetuates existing inequalities. The bias audit in
Section~\ref{sec:biasaudit} demonstrates that this layer is
not supplementary. It is the governance layer that determines
distributional outcomes independently of algorithmic quality.

\subsection{Four Acceleration Pathways}

The four layers generate four mutually reinforcing pathways
(Figure~\ref{fig:pathways}). Improved loss detection strengthens
cash flow. Stronger cash flow funds timely maintenance. Improved
maintenance enhances reliability. Greater reliability builds
willingness to pay. Improved revenue enables further data
investment, completing the loop. This is what gives the
framework its catalytic quality---gains in one domain unlock
capacity in others, but only if Layer~4 governance ensures
those gains reach all districts, not just data-rich ones.

\begin{figure}[H]
\centering
\includegraphics[width=0.90\linewidth]{DiagC_Pathways.png}
\caption{\textbf{AIBD-Cat-SL: four mutually reinforcing
acceleration pathways.}
Each pathway---reliability, financial recovery, renewable
integration, electrification planning---feeds the others
through recursive feedback. The AIBD-Cat-SL hub at the
centre coordinates all four. Without Layer~4 justice
constraints, the feedback cycle amplifies gains in
already-connected areas while bypassing the rest.}
\label{fig:pathways}
\end{figure}

%% =============================================================
\section{Quantitative Framework: Indexes and Scores}
\label{sec:math}
%% =============================================================

\subsection{Why Formal Indexes Matter}

Before introducing the three metrics that anchor this paper's
empirical contributions, it is worth explaining why we formalise
them as equations rather than describing them narratively.
The reason is accountability. Narrative descriptions of
``data poverty'' or ``algorithmic exclusion'' are hard to
audit, replicate, or contest. A numbered equation with
documented weights and published data sources is auditable.
Policymakers and procurement agencies can run the calculation
themselves with updated figures. Vendors can be required to
demonstrate their systems' performance against it. That
auditability is the entire point.

%% ── Metric 1: DRI ─────────────────────────────────────────────
\begin{tcolorbox}[metDRI,
  title={Metric 1: Data Readiness Index (DRI) ---
         Is this district ready for AI?}]

The DRI answers a practical question: can an AI system
deployed in this district be trusted to give reliable outputs?
The answer depends on whether the data environment is
sufficiently complete, consistent, and current to train
and operate a model responsibly.

For district $d$, the DRI is a PCA-weighted composite of six
observable data quality dimensions $\phi_{d,i} \in [0,100]$:

\begin{equation}
\mathrm{DRI}_d \;=\; \sum_{i=1}^{6}\, w_i \cdot \phi_{d,i},
\qquad \sum_{i=1}^{6} w_i = 1
\label{eq:dri}
\end{equation}

\noindent The six dimensions and their roles:

\medskip
\begin{tabularx}{\linewidth}{@{}lX@{}}
$\phi_1$: Metering coverage (\%) &
  Fraction of connections with functional meters \\
$\phi_2$: Billing completeness (\%) &
  Share of active accounts with current billing records \\
$\phi_3$: Outage logging (\%) &
  Share of outage events captured in structured records \\
$\phi_4$: GIS asset records (\%) &
  Fraction of distribution assets geo-referenced \\
$\phi_5$: Data recency (months) &
  Recency of last full database reconciliation \\
$\phi_6$: Interoperability (0--100) &
  Degree of cross-system data exchange standardisation \\
\end{tabularx}

\medskip
Weights $w_i$ are derived from PCA on Sierra Leone's 14-district
dataset. The first principal component explains 100\% of
composite variance, confirming that data quality is a
unidimensional structural condition. All six dimensions
move together: you cannot patch one while the others fail.

\medskip
\textbf{Critical threshold:} DRI~$< 25$ signals a district
where AI deployment is premature---outputs would be unreliable
regardless of algorithmic sophistication.
\textit{Applied:} 7 of 14 districts fall below this threshold.

\end{tcolorbox}

\bigskip

%% ── Metric 2: DBI ─────────────────────────────────────────────
\begin{tcolorbox}[metDBI,
  title={Metric 2: Double Blindspot Index (DBI) ---
         How invisible is this district to AI systems?}]

The DBI formalises the feedback loop above as a measurable,
district-level score. It captures two compounding
disadvantages simultaneously: physical deprivation (no
electricity, high poverty) and digital invisibility (data
systems cannot see you). A district with a high DBI is at
the greatest risk of algorithmic exclusion because it has
the most to gain from AI-assisted planning and the least
chance of being seen by it.

For district $d$ with electricity access rate $A_d$\,(\%),
poverty headcount $P_d$\,(\%), and data legibility score
$L_d \in [0,100]$:

\begin{equation}
\mathrm{DBI}_d \;=\;
\underbrace{\alpha_1\!\left(1 - \dfrac{A_d}{100}\right)}_{%
  \text{access deprivation}}
\;+\;
\underbrace{\alpha_2\!\left(\dfrac{P_d}{100}\right)}_{%
  \text{poverty burden}}
\;+\;
\underbrace{\alpha_3\!\left(1 - \dfrac{L_d}{100}\right)}_{%
  \text{data opacity}}
\label{eq:dbi}
\end{equation}

\noindent with weights $\alpha_1 = 0.35$,
$\alpha_2 = 0.35$, $\alpha_3 = 0.30$, and
$\mathrm{DBI}_d \in [0,100]$.

\medskip
\textbf{Weight justification:} Access deprivation and poverty
burden are weighted equally (0.35 each) because both directly
measure the need for AI-assisted planning. Data opacity (0.30)
is slightly lower because it captures the \textit{mechanism}
of exclusion rather than its direct consequence. Sensitivity
analysis confirms the core finding ($r = 0.972$) is robust
to weight variation of $\pm0.10$.

\medskip
\textbf{Applied:} DBI ranges from 15.8 in Freetown
(Western Area Urban) to 88.6 in Falaba---a 72.8-point gap
reflecting extreme spatial inequality. Falaba: 9\%
electricity access, 85\% poverty, DRI $= 0$.
Freetown: 83\% access, 18\% poverty, DRI $= 100$.

\end{tcolorbox}

\bigskip

%% ── Metric 3: Parity Gap ──────────────────────────────────────
\begin{tcolorbox}[metDP,
  title={Metric 3: Demographic Parity Gap ($\Delta_{\mathrm{DP}}$)
         --- Is the AI planning model fair?}]

The demographic parity gap operationalises the fairness
question in a form that is auditable and contractually
enforceable. For AI allocation model $\hat{Y}$ and poverty
group $G$ ($G=0$: low-poverty; $G=1$: high-poverty):

\begin{equation}
\Delta_{\mathrm{DP}} \;=\;
\bigl|\,\mathbb{P}(\hat{Y}=1 \mid G=0)
\;-\;
\mathbb{P}(\hat{Y}=1 \mid G=1)\,\bigr|
\label{eq:dp}
\end{equation}

\noindent where $\hat{Y}=1$ denotes grid extension allocation.
The fairness threshold is $\delta = 0.05$ \cite{Barocas2019}:
any AI planning system producing $\Delta_{\mathrm{DP}} > 0.05$
should not be procured without explicit equity constraints.

\medskip
\textbf{What the threshold means in practice:} If the AI
allocates grid investment to 40\% of low-poverty settlements,
it must allocate to at least 35\% of high-poverty settlements
to pass. A parity gap of 0.406 means the system gives 40\%
probability to low-poverty areas and 0\% to high-poverty
ones. That is not a marginal fairness concern. It is a
complete exclusion of the most deprived communities.

\medskip
\textbf{Applied:} Unregulated AI: $\Delta_{\mathrm{DP}} = 0.406$
(\textit{FAIL}). Justice-constrained AI:
$\Delta_{\mathrm{DP}} = 0.034$ (\textit{NEAR-PASS}).
Reduction: \textbf{92\%}.

\end{tcolorbox}

%% =============================================================
\section{Analytical Results}
\label{sec:results}
%% =============================================================

\subsection{The Double Blindspot: Spatial Evidence}

Figures~\ref{fig:spatial} and \ref{fig:blindspot} map the
DBI co-location across all 14 districts using GeoPandas
with GADM v4.1 administrative boundaries. The Pearson
correlation between electricity access rate and data
legibility across districts is $r = 0.972$.

\begin{figure}[H]
\centering
\includegraphics[width=\linewidth]{Fig3_Choropleth_Map.png}
\caption{\textbf{The Double Blindspot: spatial co-location
of physical energy poverty and digital invisibility in Sierra Leone.}
Choropleth colour encodes the DBI
(Equation~\ref{eq:dbi}: $\alpha_1=0.35$, $\alpha_2=0.35$,
$\alpha_3=0.30$). Darker red indicates districts that are
simultaneously more energy-deprived and more invisible to AI
systems. Kono/Koidu (east) and Pujehun (south) represent
the deepest double blindspot. Freetown (Western Area Urban,
west) is the lightest: most connected and most data-visible.
DBI spread: 15.8--88.6 (72.8-point gap). $r = 0.972$.
\textit{Sources:} GADM v4.1 \cite{GADM};
ESMAP/MTENR (2023) \cite{ESMAP2023};
World Bank (2021) \cite{WorldBankWDI};
DHS SL (2019) \cite{DHS2019}.}
\label{fig:spatial}
\end{figure}

\begin{figure}[H]
\centering
\includegraphics[width=\linewidth]{Fig2_Double_Blindspot.png}
\caption{\textbf{District-level Double Blindspot analysis:
scatter and DBI rankings.}
Left panel: data legibility vs.\ electricity access
(bubble colour encodes DBI). The strong upward slope confirms
the feedback loop: districts with more electricity have
more data, which means AI will see them better and direct
more investment toward them. Right panel: top 8 highest-DBI
districts by ranking---all face a ``critical'' double
blindspot of simultaneously high deprivation and low
data visibility.
\textit{Sources:} ESMAP/MTENR (2023) \cite{ESMAP2023};
World Bank (2021) \cite{WorldBankWDI};
DHS SL (2019) \cite{DHS2019};
Stats SL (2021) \cite{StatsSL2021}.}
\label{fig:blindspot}
\end{figure}

\subsection{Electrification Forecast and the Cost of
Data Gaps}

Figure~\ref{fig:forecast} presents an electrification
trajectory forecast using a polynomial trend model (degree 2)
fitted on World Bank WDI data from 1990 to 2023. Sierra Leone
is projected to reach 53.5\% national electricity access by
2031---marginally above the SDG7 50\% interim target. But the
critical output is not the projection itself. It is the right
panel: forecast uncertainty growing with each year of planning
horizon.

\begin{figure}[H]
\centering
\includegraphics[width=\linewidth]{Fig4_Forecast.png}
\caption{\textbf{Electrification trajectory forecast
(Propositions~1 and~2).}
Panel (a): polynomial trend projects 53.5\% access by 2031,
with 85\% and 95\% confidence bands. Panel (b): forecast
confidence interval width grows from 4.2\,pp at Year~1 to
4.6\,pp by Year~8. Each year without real-time data feeds
adds $\approx$0.05\,pp of compounding planning uncertainty.
This is not a statistical inconvenience. It is the
quantifiable cost of weak data infrastructure translating
into capital misallocation in a sector where every
misallocated dollar means another community left waiting.
\textit{Source:} World Bank WDI \cite{WorldBankWDI};
residual $\hat\sigma = 1.05$\,pp.}
\label{fig:forecast}
\end{figure}

\subsection{Technology Matching: AI Knowing Where to Invest}

When an AI system has access to good settlement data, it can
do something genuinely useful: tell planners whether a
community is better served by grid extension, a mini-grid,
or standalone solar. That distinction matters enormously.
Grid connections run about \$2,800 per household. Standalone
solar runs about \$220. Correct technology matching in a
country where every dollar of electricity investment is
contested is not a secondary concern. It is one of the
highest-leverage AI applications available.

Figure~\ref{fig:techmatching} shows a Random Forest classifier
achieving $F_1\text{-macro} = 0.993 \pm 0.009$ replicating
the SEforALL Global Electrification Platform logic
\cite{Korkovelos2019}. Distance to grid explains 53.6\% of
Gini feature importance---which is also the mechanism the
bias audit exposes, as discussed next.

\begin{figure}[H]
\centering
\includegraphics[width=\linewidth]{Fig5_Technology_Matching.png}
\caption{\textbf{AI technology matching matrix---Sierra Leone
(Proposition~6).}
Panel (a): Random Forest decision surface with three
viability zones. Panel (b): Gini feature importance---
distance to grid explains 53.6\%, population density 46.4\%.
Panel (c): technology distribution and per-connection cost
implications (\$2,800 grid; \$680 mini-grid; \$220 SHS).
This is the alignment problem: an AI optimising on distance
will, by construction, prefer low-poverty (near-grid)
districts, exactly as the bias audit confirms.
\textit{Sources:} ESMAP (2023) \cite{ESMAP2023};
SEforALL GEP \cite{Korkovelos2019}.}
\label{fig:techmatching}
\end{figure}

\subsection{Non-Technical Loss Detection: Revenue in Plain Sight}

Sierra Leone's EDSA loses between 15\% and 25\% of purchased
electricity to non-technical causes---meter tampering, bypass
connections, and billing irregularities \cite{ESMAP2023,
AfDB2022}. In a utility already running on thin margins,
this is not a peripheral problem. An Isolation Forest
detector calibrated to ESMAP's published NTL rates identifies
meter-bypass signatures from feeder telemetry. The physical
signal is unambiguous: bypass produces anomalously low feeder
load combined with disproportionately high voltage
drop---a violation of Ohm's Law under normal conditions
(Figure~\ref{fig:ntl}).

\begin{figure}[H]
\centering
\includegraphics[width=\linewidth]{Fig6_NTL_Detection.png}
\caption{\textbf{Non-technical loss detection
(Proposition~4).}
Panel (a): flagged customers cluster in the theft/bypass
signature zone (low load $+$ high voltage drop). Panel (b):
anomaly score distributions show clear separability.
Panel (c): cumulative 36-month revenue recovery
projection---\$78,000 USD for 60 flagged customers.
Precision/recall/$F_1 = 0.717$. Precision matters
institutionally: incorrect disconnections erode public
trust, as a poorly managed enforcement programme in Lusaka
demonstrated in 2019 when customer backlash set that
utility's digitalisation effort back two years.
\textit{Sources:} ESMAP/AfDB (2022)
\cite{ESMAP2023, AfDB2022}.}
\label{fig:ntl}
\end{figure}

\subsection{Algorithmic Bias Audit: The Paper's Centrepiece
Finding}
\label{sec:biasaudit}

Table~\ref{tab:parity} and Figures~\ref{fig:biasaudit}
and \ref{fig:justice} present the paper's most original
empirical contribution---the first demographic parity audit
of AI-driven electrification planning in an LMIC context.

\begin{table}[H]
\centering
\caption{\textbf{Demographic parity audit: grid allocation
by poverty group (Proposition~8).}
High-poverty $\equiv$ $\geq 60\%$ headcount below
\$2.15/day PPP (World Bank 2021).
Fairness threshold $\delta = 0.05$
\cite{Barocas2019}.}
\label{tab:parity}
\setlength\extrarowheight{4pt}
\begin{tabularx}{\textwidth}{|>{\arraybackslash}p{4.0cm}
                              |>{\centering\arraybackslash}p{3.3cm}
                              |>{\centering\arraybackslash}p{3.3cm}
                              |>{\centering\arraybackslash}X|}
\hline
\Th{District group} &
\Th{Unregulated AI} &
\Th{Justice-constrained} &
\Th{Change} \\
\hline
\rowcolor{rowA}
Low-poverty ($G=0$) &
  $\mathbb{P}(\hat{Y}=1) = 0.406$ &
  $\mathbb{P}(\hat{Y}=1) = 0.406$ &
  No change \\
\hline
High-poverty ($G=1$) &
  \Fail{$\mathbb{P}(\hat{Y}=1) = 0.000$} &
  \Pass{$\mathbb{P}(\hat{Y}=1) = 0.372$} &
  $+37.2$\,pp \\
\hline
\rowcolor{midGray}
\textbf{Parity gap} $\Delta_{\mathrm{DP}}$ &
  \Fail{0.406 \textemdash FAIL} &
  \Pass{0.034 \textemdash NEAR-PASS} &
  $\boldsymbol{-92\%}$ \\
\hline
\multicolumn{4}{|l|}{\small
  Poverty--selection correlation (unregulated AI,
  14 districts): $r = -0.754$}\\
\hline
\end{tabularx}
\end{table}

\begin{figure}[H]
\centering
\includegraphics[width=\linewidth]{Fig7_Algorithmic_Bias_Audit.png}
\caption{\textbf{Algorithmic bias audit: grid allocation
fairness (Proposition~8---centrepiece finding).}
Panel (a): unregulated AI gives high-poverty districts
zero allocation probability ($\Delta_{\mathrm{DP}} = 0.406$);
justice-constrained design reduces the gap by 92\% to
$\Delta_{\mathrm{DP}} = 0.034$.
Panel (b): poverty--selection scatter across 14 districts,
$r = -0.754$.
Panel (c): parity gap reduction.
Panel (d): intersectional analysis (poverty $\times$
distance to grid)---justice constraint is most impactful
at 25--150\,km, exactly where the most deprived
districts sit.
Panel (e): key findings summary.
\textit{Sources:} World Bank \cite{WorldBankWDI};
ESMAP \cite{ESMAP2023}; Stats SL \cite{StatsSL2021};
Barocas et al.\ \cite{Barocas2019}.}
\label{fig:biasaudit}
\end{figure}

\begin{figure}[H]
\centering
\includegraphics[width=\linewidth]{Fig7_Algorithmic_Bias_Audit.png}
\caption{\textbf{Algorithmic justice: mitigating spatial
redlining in electricity allocation.}
Urban (low-poverty) districts receive 88\% acceptance
rate under unregulated AI; rural (high-poverty) districts
receive just 13\%. Justice-constrained AI
(AIBD-Cat-SL Layer~4) reduces the urban acceptance rate
slightly to 55\% while raising rural districts to 45\%---
a near-equitable allocation without sacrificing system
efficiency.
\textit{Source:} Authors' analysis; ESMAP (2023)
\cite{ESMAP2023}.}
\label{fig:justice}
\end{figure}

The mechanism is structural, not intentional. Every high-poverty
district in Sierra Leone is located more than 90\,km from
the existing grid---a geographic reality created by decades
of unequal investment. An AI optimising on distance will
give all of them zero allocation probability. It is not
making a mistake. It is doing exactly what it was designed
to do---and in doing so, it encodes historical inequality
as future infrastructure allocation.

The justice constraint is a single governance rule: extend
coverage eligibility to 60\,km for settlements where the
poverty proxy exceeds 70\%. This raises high-poverty
district allocation rates to 0.372, reduces the parity
gap by 92\%, and leaves Freetown's allocation rate
unchanged. No algorithm change. No additional data collection.
One procurement requirement.

\subsection{Data Readiness Index: Where AI Can and Cannot
Work}
\label{sec:dri}

The DRI (Equation~\ref{eq:dri}) is a PCA-weighted composite
of six data quality dimensions from published ESMAP and EDSA
reports \cite{ESMAP2023, EDSA2022} (Figure~\ref{fig:dri}).

\begin{figure}[H]
\centering
\includegraphics[width=\linewidth]{Fig8_Data_Readiness_Index.png}
\caption{\textbf{Data Readiness Index by district
(Proposition~1 operationalised).}
Panel (a): DRI scores by district with tier classifications.
Seven districts fall in the Critical tier (DRI~$< 25$)
where AI deployment is premature. Panel (b): dimension
comparison between Western Area Urban (DRI~$= 100$) and
Falaba (DRI~$= 0$)---a total structural gap across all
six dimensions. Panel (c): DRI vs.\ access scatter,
$r = 0.97$, confirming that data readiness and
electricity access are essentially the same variable in
two different units. The Double Blindspot Zone in panel
(c) marks where AI risk is greatest.
\textit{Sources:} ESMAP (2023) \cite{ESMAP2023};
EDSA (2022) \cite{EDSA2022}; MTENR (2022)
\cite{MTENR2022}.}
\label{fig:dri}
\end{figure}

Table~\ref{tab:results} consolidates all key quantitative
findings.

\begin{table}[H]
\centering
\caption{\textbf{Summary of quantitative findings.}
The highlighted row (purple) marks the centrepiece empirical
contribution.}
\label{tab:results}
\setlength\extrarowheight{3pt}
\begin{tabularx}{\textwidth}{|>{\centering\arraybackslash}p{0.9cm}
                              |>{\arraybackslash}p{2.8cm}
                              |>{\arraybackslash}X
                              |>{\centering\arraybackslash}p{1.3cm}|}
\hline
\Th{Fig.} & \Th{Demonstration} &
\Th{Key quantitative result} & \Th{Prop.} \\
\hline
\rowcolor{rowA}
1 & Macro trends &
  Access 7.7\% (2000) $\to$ 35.5\% (2023);
  losses 30\% ($+10$\,pp above benchmark) & --- \\
\hline
2, 3 & Double Blindspot &
  DBI 15.8--88.6; $r = 0.972$; Falaba worst (88.6) &
  P1, P8 \\
\hline
\rowcolor{rowA}
4 & Forecast &
  53.5\% by 2031; CI widens 4.2--4.6\,pp;
  $\hat\sigma = 1.05$\,pp & P1, P2 \\
\hline
5 & Technology matching &
  RF $F_1 = 0.993 \pm 0.009$;
  distance explains 53.6\% Gini importance & P6 \\
\hline
\rowcolor{rowA}
6 & NTL detection &
  $F_1 = 0.717$; 60 customers flagged;
  \$78K revenue recovery (36\,mo) & P4 \\
\hline
\rowcolor{rowD}
\textbf{7} & \textbf{Bias audit} &
  $\boldsymbol{\Delta_{\mathrm{DP}} = 0.406 \to 0.034}$;
  \textbf{92\% reduction}; $r = -0.754$ &
  \textbf{P8} \\
\hline
\rowcolor{rowA}
8 & Data Readiness Index &
  $r(\text{DRI, access}) = 0.97$;
  7 of 14 districts Critical (DRI~$< 25$) & P1 \\
\hline
\end{tabularx}
\end{table}

%% =============================================================
\section{Monte Carlo Validation: From Point Estimates to
Structural Properties}
\label{sec:mc}
%% =============================================================

\subsection{Why Monte Carlo?}

Every finding in Section~\ref{sec:results} is a point estimate
derived from calibrated published data. A reviewer is entitled
to ask: what if the district statistics are slightly wrong?
What if the DBI weights were calibrated differently? Is the
92\% gap reduction a one-run finding, or a genuine structural
property of the data environment?

Monte Carlo simulation answers this directly. By running
thousands of independent realisations with controlled noise
applied to each indicator, we test whether the findings are
robust to measurement uncertainty or sensitive to specific
parameter choices. This validation is especially important
given the paper's core argument: that Sierra Leone faces a
\textit{structural} algorithmic deficit rooted in the
geographic co-location of physical and digital deprivation.
Structural arguments require structural validation.

\begin{tcolorbox}[metMC,
  title={Monte Carlo Bootstrap Estimator}]

For each bootstrap iteration $b = 1, \ldots, B$, indicator
$k$ (access rate, poverty headcount, data legibility) for
district $d$ is perturbed with $\pm 5\%$ Gaussian noise:

\begin{equation}
x_{d,k}^{(b)} \;=\; x_{d,k}\cdot
\bigl(1 + \varepsilon_{d,k}^{(b)}\bigr),
\qquad
\varepsilon_{d,k}^{(b)} \sim \mathcal{N}(0,\,\sigma^2),
\quad \sigma = 0.05
\label{eq:bootstrap}
\end{equation}

The bootstrap estimator and 95\% confidence interval for
any district-level statistic $\theta$ (DBI rank, parity gap,
or revenue recovery):

\begin{equation}
\hat\theta_{\mathrm{boot}} \;=\;
\tfrac{1}{B}\sum_{b=1}^{B}\theta^{(b)},
\qquad
\mathrm{CI}_{95\%} =
\bigl[\hat\theta_{(0.025)},\;\hat\theta_{(0.975)}\bigr]
\label{eq:bootci}
\end{equation}

$B = 10{,}000$ for DBI robustness; $B = 5{,}000$ for parity
gap; $B = 10{,}000$ for NTL revenue recovery. All simulations
use real published district statistics as the baseline.
The $\sigma = 0.05$ perturbation represents a conservative
$\pm 5\%$ bound on administrative statistics in data-scarce
settings---if anything, the real uncertainty may be larger,
which makes the robustness findings more striking.

\end{tcolorbox}

\subsection{MC~1: Is Falaba Always Worst?}

Across 10,000 bootstrap iterations, Falaba ranks among the
three worst districts in \textbf{99.9\% of runs}.
Western Area Urban never appears in the worst three. The
95\% confidence intervals on DBI scores are narrow
(half-width: 2--4 points), and the DBI distributions for
Falaba and Freetown show zero overlap across all 10,000
simulations (Figure~\ref{fig:mc_dbi}).

\begin{figure}[H]
\centering
\includegraphics[width=\linewidth]{Fig9_MC_DBI_Robustness.png}
\caption{\textbf{Bootstrap DBI robustness: 10,000 simulations
(Equations~\ref{eq:bootstrap}--\ref{eq:bootci}).}
Panel A: mean DBI with 95\% bootstrap CIs---narrow bands
confirm stable estimates even under data noise. Panel B:
rank stability---Falaba is ranked among the worst 3 districts
in 99.9\% of runs; Koinadugu reaches 92.7\%. Panel C:
the DBI distributions for Falaba (mean 88.6) and Western
Area Urban (mean 15.8) have zero overlap across all 10,000
simulations. The Double Blindspot is not a measurement
artefact.}
\label{fig:mc_dbi}
\end{figure}

\subsection{MC~2: Is the Parity Gap Always Above the
Fairness Threshold?}

Across 5,000 independent Monte Carlo runs generating
settlement populations from real district data with randomised
distance noise (log-normal, $\sigma = 0.55$) and poverty
proxy noise (Gaussian, $\sigma = 0.05$)
(Figure~\ref{fig:mc_parity}):

\begin{itemize}[leftmargin=1.5em, itemsep=2pt]
\item The \textbf{unregulated AI parity gap exceeds
  $\delta = 0.05$ in 100\% of runs}---no parameterisation
  of Sierra Leone's geography produces a fairness-passing
  cost-optimised AI.
\item The \textbf{justice constraint reduces the gap in
  99\% of runs}---the directional effect is robust to
  parameter variation.
\end{itemize}

\begin{figure}[H]
\centering
\includegraphics[width=\linewidth]{Fig10_MC_Parity_Gap.png}
\caption{\textbf{Monte Carlo parity gap distribution:
5,000 runs.}
Panel A: unregulated AI (red) exceeds the 0.05 fairness
threshold in 100\% of simulations; justice-constrained
AI (green) reduces the gap in 99\% of runs---both
distributions clearly separated from the threshold.
Panel B: gap reduction distribution, mean 3.9\%, 5th
percentile 0.9\%---positive in virtually every run.
Panel C: selection rate scatter (400 of 5,000 runs
shown)---unregulated cloud clusters near zero for
high-poverty; justice cloud moves upward toward the
parity line.}
\label{fig:mc_parity}
\end{figure}

\subsection{MC~3: What Is the Realistic Revenue Recovery Range?}

Sampling jointly from ESMAP/AfDB published ranges for NTL
rate (Uniform 15--25\%), per-customer monthly loss (Uniform
USD 28--45), and detection precision (Beta(7,3), mean
$\approx 0.70$) across 10,000 runs (Figure~\ref{fig:mc_ntl}):

\begin{itemize}[leftmargin=1.5em, itemsep=2pt]
\item \textbf{Mean 36-month revenue recovery: USD 67,000}
  (90\% credible interval: USD 37,000--101,000).
\item \textbf{P(recovery $>$ USD 60K) = 61\%}---more likely
  than not even at pessimistic detection rates.
\end{itemize}

\begin{figure}[H]
\centering
\includegraphics[width=\linewidth]{Fig11_MC_NTL_Revenue.png}
\caption{\textbf{Monte Carlo NTL revenue recovery: 10,000
simulations.}
Panel A: 36-month recovery distribution, mean USD 67K,
90\% CI [USD 37K--101K]. Panel B: recovery vs.\ NTL rate
(colour = detection precision)---positive linear
relationship confirms that investment in metering
infrastructure that reduces NTL uncertainty has
compounding financial returns. Panel C: cumulative
exceedance probability---P(recovery~$>$~USD 50K) = 80\%.
\textit{Sources:} ESMAP (2023) \cite{ESMAP2023};
AfDB (2022) \cite{AfDB2022}.}
\label{fig:mc_ntl}
\end{figure}

%% =============================================================
\section{Research Propositions}
\label{sec:props}
%% =============================================================

\begin{table}[H]
\centering
\caption{\textbf{Eight empirically testable propositions.}
Propositions with current analytical support are marked
Supported; those requiring future field data are marked
Requires fieldwork.}
\label{tab:props}
\setlength\extrarowheight{3pt}
\begin{tabularx}{\textwidth}{|>{\centering\arraybackslash}p{0.6cm}
                              |>{\arraybackslash}X
                              |>{\arraybackslash}p{3.6cm}|}
\hline
\ThP{P} & \ThP{Proposition} & \ThP{Evidence status} \\
\hline
\rowcolor{rowA}
P1 & AI adoption outcomes are moderated by data readiness;
     AI produces more reliable outputs where quality,
     interoperability, and temporal continuity are
     stronger \cite{Ahmad2021, Kitchin2014}. &
  \textit{Supported}: DRI $r=0.97$; 7/14 Critical;
  MC bootstrap; CI growth \\
\hline
P2 & Institutional capacity mediates the analytics--performance
     link; accurate models fail where organisations
     lack routines to act on outputs \cite{Danish2023}. &
  \textit{Supported}: forecast CI growth horizon \\
\hline
\rowcolor{rowA}
P3 & Explainability and procedural transparency increase
     adoption and sustained use of AI in high-stakes
     energy decisions \cite{Machlev2022}. &
  \textit{Requires} field adoption studies \\
\hline
P4 & AI-enabled loss analytics strengthen utility viability
     when combined with field verification, enforcement
     protocols, and appeals mechanisms
     \cite{Chatterjee2022}. &
  \textit{Supported}: NTL demo;
  MC recovery USD 37K--101K \\
\hline
\rowcolor{rowA}
P5 & Predictive maintenance yields larger welfare effects
     in weak-grid environments where avoided downtime
     produces proportionally greater gains. &
  \textit{Requires} pre/post field comparison \\
\hline
P6 & Geospatial AI reduces capital misallocation by improving
     technology matching \cite{Korkovelos2019, Arevalo2024}. &
  \textit{Supported}: RF $F_1=0.993$;
  \$2,800 vs \$220 cost gap \\
\hline
\rowcolor{rowA}
P7 & Cybersecurity maturity is a necessary enabling condition
     for scaled digitalisation; inadequate safeguards
     heighten systemic risk. &
  \textit{Requires} infrastructure security audit \\
\hline
\rowcolor{rowD}
\textbf{P8} &
  \textbf{The developmental value of AI is greatest when
  aligned with energy justice; unregulated optimisation
  reproduces spatial inequality; justice-constrained
  design closes the gap}
  \cite{Sovacool2015, Barocas2019}. &
  \textbf{Supported}: $\Delta_{\mathrm{DP}}=0.406\to0.034$;
  92\% reduction; MC 99\% of 5,000 runs \\
\hline
\end{tabularx}
\end{table}

%% =============================================================
\section{Implications for Empirical Research Design}
\label{sec:design}
%% =============================================================

A persistent weakness in this literature is calling for
further empirical research in broad terms without specifying
what that research should look like. This section avoids
that trap.

\textit{Feeder-level panel analysis.} Using EDSA operational
data, estimate whether AI-enabled loss analytics improve
commercial recovery, outage frequency, and service quality
across feeders with different levels of data maturity.
This directly tests Proposition~1: does higher DRI produce
better AI outcomes?

\textit{Mini-grid portfolio studies.} Use telemetry data and
maintenance logs to examine whether remote monitoring reduces
downtime and lowers lifecycle costs in Sierra Leone's growing
off-grid sector. This tests Proposition~5, where the MC
analysis predicts the largest welfare gains.

\textit{Technology matching verification.} Compare AI-informed
siting decisions with conventional planning approaches
in terms of cost per connection, time to service delivery,
and reliability. This validates Proposition~6 with primary
implementation evidence rather than simulated data.

\textit{Institutional ethnography.} Many AI failures run
through procurement choices, staffing arrangements, and the
politics of reform---not through algorithmic quality.
Process tracing and organisational ethnography are
indispensable for explaining why interventions scale or
stall. Particularly important: how is the parity gap
requirement embedded in procurement contracts, and who
enforces it?

The most powerful research designs will combine quantitative
outcome estimation with institutional mechanism tracing.
Randomised evaluation, if feasible, would be ideal for
the NTL detection programme. Difference-in-differences
designs could be used for feeder-level analytics.
Qualitative case comparison is appropriate for the
institutional adoption questions.

%% =============================================================
\section{Policy and Implementation Roadmap}
\label{sec:roadmap}
%% =============================================================

The governing principle is \textit{informational order before
computational sophistication}. The DRI analysis shows that
seven of 14 districts are in a Critical tier where AI outputs
would be unreliable regardless of algorithmic quality.
No amount of machine learning investment changes that.
The data foundation must come first.

Figure~\ref{fig:roadmap} presents the three-phase roadmap.

\begin{figure}[H]
\centering
\includegraphics[width=\linewidth]{DiagB_Roadmap.png}
\caption{\textbf{AIBD-Cat-SL: three-phase implementation
roadmap.}
Phase~1 (Years 1--2): data foundation and NTL pilots.
The six Critical-DRI districts require data infrastructure
investment before AI can be responsibly deployed. This is
what Ghana did between 2008 and 2012 before scaling digital
metering sector-wide. Phase~2 (Years 3--5): mandatory parity
auditing as a procurement condition, operationalising
Proposition~8 as policy. Phase~3 (Year~6+): system
intelligence and full algorithmic accountability; success
measured by equity improvement, not tool sophistication.}
\label{fig:roadmap}
\end{figure}

\paragraph{Phase~1 (Years 1--2).}
National energy data standards. Comprehensive asset
registries. Metering-billing integration. Cybersecurity
baseline. Feeder-level NTL pilots with transparent,
auditable outputs and explicit parity monitoring. The DBI
analysis ($r = 0.972$) establishes the spatial priority
order for data investment: start with the seven Critical-DRI
districts and work outward.

\paragraph{Phase~2 (Years 3--5).}
Mandatory demographic parity reporting as a condition
of AI procurement approval. Any AI planning system
deployed by EDSA or procured through development partner
funding should be required to demonstrate
$\Delta_{\mathrm{DP}} < 0.05$ on Sierra Leone district
data before deployment. The bias audit result
($\Delta_{\mathrm{DP}} = 0.406$ under unregulated AI)
provides the empirical baseline against which systems
must be benchmarked. Analytics must be embedded in
maintenance schedules, inspection routines, investment
decisions, and regulatory review---not just dashboard
displays.

\paragraph{Phase~3 (Year~6+).}
Digital twins, integrated geospatial platforms for
on/off-grid coordination, formal algorithmic accountability
mechanisms. The measure of success is not the sophistication
of tools deployed but whether the parity gap closes,
outages decline, and communities in Falaba and Koinadugu
receive the same quality of planning attention as
communities in Freetown.

%% =============================================================
\section{Discussion}
\label{sec:discussion}
%% =============================================================

\subsection{Three Deficits, Not Two}

The central theoretical contribution of this article is to
identify and operationalise an algorithmic deficit alongside
the familiar infrastructural and financial ones. Most of the
existing literature implicitly treats digital transformation
as a subsequent modernisation layer added to already-functional
systems \cite{Arevalo2024, Yousef2023}. Sierra Leone invites
a different reading: in fragile, only partially legible
systems, digital capability is not an accessory to transition.
It is part of the transition's foundational challenge. Where
the system cannot see itself, it cannot improve itself.

The DBI finding ($r = 0.972$) extends
\citeauthor{Kitchin2014}'s \citeyearpar{Kitchin2014}
data-as-infrastructure argument by demonstrating its spatial
dimension: the physical deficit and the data deficit
co-locate with extraordinary precision. This has a
consequential implication for development sequencing.
A district that lacks electricity access also lacks the data
infrastructure that would enable an AI system to allocate
it investment. Breaking this cycle requires simultaneous
investment in physical connectivity and data infrastructure---
not sequential priority.

\subsection{Governance as the Decisive Variable}

The bias audit moves the energy justice debate
\cite{Sovacool2015} from normative aspiration to empirical
accountability. A parity gap of 0.406 is not a rhetorical
claim about inequality. It is a measurable, auditable,
reducible property of an allocation system. The justice-
constrained design achieves a 92\% reduction without
sacrificing efficiency. This establishes a concrete
procurement specification: parity gap reporting, maximum
gap thresholds ($\Delta_{\mathrm{DP}} < 0.05$), and audit
cycles as standard procurement conditions, analogous to
the financial and technical performance standards already
routinely required of infrastructure contractors.

What makes this finding particularly significant is that
the 92\% reduction came not from a more sophisticated
algorithm but from one governance rule. The algorithm is
fine. The governance is the problem---and it is also the
solution.

\subsection{Epistemic Reconfiguration as a Transition
Condition}

The paper contributes to sociotechnical transition theory
\cite{Geels2002} by insisting that regime transformation
requires not only material and institutional change but
\textit{epistemic reconfiguration}: new ways of seeing,
measuring, and coordinating a system growing more complex
faster than traditional modes can manage. Data standards,
telemetry architectures, and interoperable information flows
are not peripheral complements to transition---they are
components of the architecture through which transition
becomes governable in practice.

\subsection{Limitations}

Three limitations warrant acknowledgement. First, the
analytical demonstrations use calibrated rather than
directly operational data: published statistics
parameterise the models rather than primary EDSA data.
The bias audit is therefore indicative, not definitively
causal. Second, the paper focuses on electricity and
underweights clean cooking, which remains central to
Sierra Leone's broader energy transition and to gendered
welfare. The DBI and parity audit logic should be extended
to clean cooking access in future work. Third, the paper
identifies governance requirements without providing
complete institutional design---legal mandates, regulatory
enforcement mechanisms, and procurement standards require
dedicated subsequent work.

%% =============================================================
\section{Conclusion}
\label{sec:conclusion}
%% =============================================================

Sierra Leone's energy transition is most productively
understood not as a problem of building more power plants
or laying more cable, but as a challenge of constructing the
informational and governance capabilities necessary to
observe, interpret, and equitably distribute the benefits
of an increasingly complex electricity system.

Start with what the data actually shows. In Falaba:
9\% electricity access, 85\% poverty, DBI of 88.6, AI
investment probability of 0.000. In Freetown: 83\%
access, 18\% poverty, DBI of 15.8, AI investment
probability of 0.406. The parity gap is 0.406.
This is not a simulation artefact. It is the structural
output of applying a cost-optimisation system to a data
environment shaped by fifty years of unequal investment.

Change one governance rule. The parity gap falls to 0.034.
Freetown's probability does not change. No algorithm is
rewritten. No data is collected. One procurement requirement.

The AIBD-Cat-SL framework provides the architecture through
which this governance design can be co-developed with
physical infrastructure and digital capability. The
Double Blindspot Index ($r = 0.972$ across 14 districts)
establishes the spatial priority order for investment.
Monte Carlo validation across 15,000 simulations confirms
these are structural findings, not parameter-dependent
artefacts. The eight propositions and three-phase roadmap
translate the framework into sequenced, actionable guidance
grounded in the actual data conditions of Sierra Leone's 14
districts.

The broader message is direct. When AI systems are deployed
in settings shaped by historical inequality, technical
optimisation without explicit equity constraints does not
close the gap. It encodes it. The AIBD-Cat-SL framework
provides the structured basis for ensuring that digital
capability becomes a catalyst of equitable transition---not
another instrument of differential advantage dressed in
the language of innovation.

Any organisation currently procuring AI for electrification
planning in Sierra Leone or any comparable setting should
ask one question before signing a contract: \textit{what
is the demographic parity gap of this system, and what
governance constraint ensures it stays below 0.05?} If
the vendor cannot answer, the system has not been audited
for equity. This paper provides the method, the threshold,
and the precedent. The rest is governance.

\section*{Acknowledgements}
The authors thank colleagues at Thammasat University and
Carnegie Mellon University Africa for critical feedback on
earlier drafts. We acknowledge open-data contributions
from the World Bank, GADM (University of California, Davis),
the Energy Sector Management Assistance Programme, Statistics
Sierra Leone, the DHS Programme, UNDP, and the Global
Electrification Platform (SEforALL/OnSSET). All analytical code, 
high-resolution figures, and harmonized datasets are openly 
available at: \url{https://github.com/Nakamoto-ctrl/sierra-leone-energy-transition}.

\section*{Funding}
This research did not receive any specific grant from
funding agencies in the public, commercial, or
not-for-profit sectors.

\section*{Declaration of competing interest}
The authors declare no competing interests.

\section*{Data availability}
All data are drawn from open-access sources listed in
Table~\ref{tab:data}. The complete production environment, including 
Jupyter Notebooks and raw data tables, is maintained at: 
\url{https://github.com/Nakamoto-ctrl/sierra-leone-energy-transition}.

\section*{Declaration of generative AI use}
During the preparation of this work the authors used Python
with Matplotlib, scikit-learn, NumPy, SciPy, and GeoPandas
to produce all figures and statistical analyses. These are
deterministic, reproducible programmatic libraries, not
generative AI services. No large language model or image
generation model was used to create, alter, or draft any
figure, image, or textual content in this manuscript.
The authors reviewed and take full responsibility for all content.

\bibliographystyle{elsarticle-harv}
\bibliography{references}

\end{document}